\section{Introducción}
Las redes constituyen la base operativa sobre la cual las 
organizaciones sostienen sus procesos de comunicación y acceso a la 
información. Mantener la visibilidad sobre lo que ocurre dentro de estas 
redes es una condición necesaria para preservar la disponibilidad de los 
servicios y la seguridad de los activos conectados. Sin embargo, las 
herramientas de monitoreo tradicionales operan desde un punto fijo 
mediante puertos mirror o interfaces observadas
~\cite{kurose2010redes}, lo que ofrece una visión agregada del tráfico 
pero carece de granularidad por dispositivo, y obliga al administrador 
a permanecer anclado a una estación de trabajo para consultar esa 
información. Cualquier desplazamiento genera un intervalo sin vigilancia 
durante el cual la red sigue operando pero nadie observa su 
comportamiento.

El presente trabajo aborda esta problemática mediante el desarrollo de 
GOATGuard, un sistema de monitoreo distribuido con acceso móvil 
compuesto por tres componentes: un sensor de captura desplegable en 
cada endpoint de la red, un backend centralizado que recibe, procesa y 
almacena la información reportada por los sensores, y una aplicación 
móvil desde la cual el administrador consulta el estado de la red y 
recibe notificaciones ante comportamientos inusuales. Los sensores 
capturan el tráfico directamente en el endpoint donde operan, lo que 
permite correlacionar cada flujo de red con el dispositivo que lo 
originó sin exponer el contenido del payload~\cite{stevens2004unix}. 
El backend delega la inspección profunda de protocolos al framework 
Zeek~\cite{paxson1999zeek} y aplica un motor de detección de anomalías 
basado en media móvil ponderada exponencialmente con Z-Score adaptativo 
~\cite{hunter1986ewma}, capaz de construir un perfil individual por 
dispositivo sin recurrir a umbrales fijos ni a modelos de aprendizaje 
automático externos.

El aporte del proyecto no reside en la introducción de una técnica 
aislada, sino en la integración de captura distribuida, procesamiento 
contextualizado, inventario dinámico de activos y acceso móvil fuera 
de una arquitectura unificada, orientada a redes de pequeña y mediana 
escala donde las alternativas comerciales resultan económicamente 
inviables u operativamente sobredimensionadas. El desarrollo articula 
conocimientos en redes de datos, seguridad informática, ingeniería de 
software y desarrollo móvil, consistente con la naturaleza integradora 
del proyecto académico-investigativo en el que se enmarca.